\chapter{Project Tracks for Student Interests}

Students should be able to bring their own interests to the program. These tracks help adapt the same uafR workflow to different fields while preserving a shared quality standard.

\section{Track 1: Plant, Agriculture, and Natural Products}
\textbf{Good questions:}
\begin{itemize}
  \item Do volatile profiles differ between treatments, cultivars, tissues, or developmental stages?
  \item Which detected compounds have natural product occurrence evidence?
  \item Do compounds cluster by sensory traits, pathway context, or functional groups?
\end{itemize}

\textbf{High-value tables:}
\begin{itemize}
  \item \texttt{ChemicalOccurrences}, \texttt{ChemicalTaxonomy}, and \texttt{LOTUSProfile}.
  \item \texttt{ChemicalPathwayRoles}, \texttt{KEGGPathways}, and \texttt{ChemicalTraits}.
\end{itemize}

\textbf{Recommended figures:} Occurrence heatmap, trait matrix heatmap, pathway role summary, compound class bar chart.

\textbf{Caution:} Natural product occurrence in another organism supports hypothesis generation, not direct evidence of occurrence in the student's sample.

\section{Track 2: Environmental Exposure and Toxicology}
\textbf{Good questions:}
\begin{itemize}
  \item Which detected chemicals have hazard, exposure, or regulatory signals requiring follow-up?
  \item Do sample groups differ by safety-relevant trait profiles?
  \item Which compounds have properties suggesting persistence, volatility, or hydrophobicity?
\end{itemize}

\textbf{High-value tables:}
\begin{itemize}
  \item \texttt{ChemicalHazards}, \texttt{SafetyProfile}, and \texttt{PubChemProperties}.
  \item \texttt{ChemicalMeasurementSummary}, \texttt{ChemicalTraits}, and \texttt{DerivedGroups}.
\end{itemize}

\textbf{Recommended figures:} Hazard code count, property scatterplot, source coverage heatmap, trait confidence matrix.

\textbf{Caution:} Hazard annotations are screening evidence. Risk requires dose, exposure, route, organism, and context.

\section{Track 3: Food, Flavor, Aroma, and Sensory Chemistry}
\textbf{Good questions:}
\begin{itemize}
  \item Which compounds have flavor, odor, or sensory descriptors relevant to a food or plant system?
  \item Do treatment groups differ in sensory-relevant chemical trait profiles?
  \item Which compounds should be prioritized for sensory validation or standards?
\end{itemize}

\textbf{High-value tables:}
\begin{itemize}
  \item \texttt{FEMAProfile}, \texttt{ChemicalTerms}, and \texttt{ChemicalTraits}.
  \item \texttt{PubChemProperties} and \texttt{ChemicalMeasurementSummary}.
\end{itemize}

\textbf{Recommended figures:} Sensory term heatmap, compound abundance versus sensory evidence, trait count by treatment.

\textbf{Caution:} A flavor annotation does not prove perceptibility. Concentration, threshold, matrix, and mixture effects matter.

\section{Track 4: Biomedical, Pharmacology, and Public Health}
\textbf{Good questions:}
\begin{itemize}
  \item Which compounds have biomedical terms, pharmacologic actions, or assay evidence?
  \item Are any detected compounds connected to targets, potencies, or bioactivity categories?
  \item Which compounds require caution before being described as therapeutic, toxic, or mechanistic?
\end{itemize}

\textbf{High-value tables:}
\begin{itemize}
  \item \texttt{MeSHProfile}, \texttt{PubChemBioactivity}, and \texttt{ChemicalTraitEvidence}.
  \item \texttt{ChemicalBioactivities}, \texttt{ChemicalTargets}, and \texttt{ChemicalPotencies}.
\end{itemize}

\textbf{Recommended figures:} Bioactivity summary table, target evidence table, biomedical trait matrix, evidence ladder plot.

\textbf{Caution:} Bioassay or pharmacologic annotation does not prove clinical relevance or mechanism in the student's system.

\section{Track 5: Methods, Data Science, and Tool Evaluation}
\textbf{Good questions:}
\begin{itemize}
  \item How do trait matrix settings change clustering or grouping conclusions?
  \item Which database sources contribute the most useful evidence for a chemical class?
  \item How robust are conclusions when low-confidence traits are removed?
\end{itemize}

\textbf{High-value tables:}
\begin{itemize}
  \item \texttt{ValidationSummary}, \texttt{TableQuality}, and \texttt{SourceDiagnostics}.
  \item \texttt{SourceCoverage}, \texttt{ChemicalTraitMatrix}, and \texttt{ChemicalTraitSimilarity}.
\end{itemize}

\textbf{Recommended figures:} Matrix sensitivity comparison, source coverage summary, trait overlap network, confidence distribution.

\textbf{Caution:} Method evaluation projects still need a scientific question. They should not become demonstrations of every available output.

\section{Track 6: Chemistry Education and Communication}
\textbf{Good questions:}
\begin{itemize}
  \item Which chemical traits are easiest or hardest for new students to interpret correctly?
  \item Can a small set of compounds teach uncertainty, source coverage, and chemical grouping?
  \item What visualizations best communicate chemical evidence to non-specialists?
\end{itemize}

\textbf{High-value tables:}
\begin{itemize}
  \item \texttt{ChemicalTraitReport} and \texttt{ChemicalTraitEvidence}.
  \item \texttt{DataDictionary} and \texttt{SourceCoverage}.
\end{itemize}

\textbf{Recommended products:} Teaching dataset, annotated figure set, evidence interpretation guide, mini-workshop.

\textbf{Caution:} Educational products should preserve scientific accuracy. Simplification should not erase uncertainty.

\section{Capstone Options}
Students can finish with one of several products:
\begin{itemize}
  \item Research report suitable for a dsDNA Core or lab meeting.
  \item Poster-ready figure set with captions and methods.
  \item Manuscript-style results section.
  \item Reusable analysis pipeline for a larger project.
  \item Database evidence atlas for a chemical class.
  \item Teaching module for future interns.
\end{itemize}

\section{Track Selection Worksheet}
\begin{tabularx}{\textwidth}{p{0.28\textwidth}Y}
\toprule
\textbf{Prompt} & \textbf{Response}\\
\midrule
Chosen track & \rule{0pt}{2.5em}\\
Why this track fits the student's interest & \rule{0pt}{3em}\\
Primary uafR tables & \rule{0pt}{3em}\\
Most important figure & \rule{0pt}{3em}\\
Most important limitation & \rule{0pt}{3em}\\
Likely final product & \rule{0pt}{3em}\\
\bottomrule
\end{tabularx}
